\documentclass[11pt, a4paper]{article}

% ── Packages ──────────────────────────────────────────────────────────
\usepackage[utf8]{inputenc}
\usepackage[T1]{fontenc}
\usepackage{lmodern}
\usepackage[margin=1in]{geometry}
\usepackage{amsmath, amssymb, amsthm}
\usepackage{booktabs}
\usepackage{longtable}
\usepackage{graphicx}
\usepackage{hyperref}
\usepackage{xcolor}
\usepackage{listings}
\usepackage{tcolorbox}
\usepackage{polya}

\hypersetup{
  colorlinks=true,
  linkcolor=polyablue,
  urlcolor=polyablue,
  citecolor=polyablue,
}

% ── Code listing style ────────────────────────────────────────────────
\lstdefinestyle{polya-python}{
  language=Python,
  basicstyle=\ttfamily\small,
  keywordstyle=\color{polyablue}\bfseries,
  commentstyle=\color{polyagray}\itshape,
  stringstyle=\color{polyagreen},
  numbers=left,
  numberstyle=\tiny\color{polyagray},
  numbersep=8pt,
  frame=single,
  framerule=0.4pt,
  rulecolor=\color{polyagray},
  breaklines=true,
  showstringspaces=false,
  tabsize=4,
}
\lstset{style=polya-python}
\title{Mortgage Comparison Analysis}
\author{uber-polya}
\date{2026-02-26}

\begin{document}

\maketitle
\tableofcontents
\newpage

% ── Phase A: Problem Formulation ─────────────────────────────────────
\section{Problem Formulation}

\begin{polyamodel}

\textbf{Problem Type}: Find \hfill
\textbf{Domain}: Financial Mathematics


\end{polyamodel}

% ── Universe ──────────────────────────────────────────────────────────
\subsection{Universe}

\begin{itemize}
  \item $Loan amount: $320,000$
  \item $Options: \{30-Year Fixed 6.5%, 15-Year Fixed 5.8%, Refinance 30yr->25yr at 5.2%\}$
\end{itemize}

% ── Variables ─────────────────────────────────────────────────────────
\subsection{Variables}

\begin{longtable}{lll}
\toprule
\textbf{Symbol} & \textbf{Meaning} & \textbf{Type / Range} \\
\midrule
\endhead
$L$ & loan principal = \$320,000 & $\mathbb{{R}}_{>0}$ \\
$r$ & monthly interest rate = APR/12 & $\mathbb{{R}}_{>0}$ \\
$n$ & total number of monthly payments & $\mathbb{{Z}}^+$ \\
$PMT$ & monthly payment amount & $\mathbb{{R}}_{>0}$ \\
\bottomrule
\end{longtable}

% ── Structure ─────────────────────────────────────────────────────────
\subsection{Structure}

A homebuyer is purchasing a \$400,000 home with an \$80,000 down payment (\$320,000 loan). Compare three mortgage options:

1. **30-year fixed** at 6.5\% APR
2. **15-year fixed** at 5.8\% APR
3. **Refinance scenario**: Start with the 30-year at 6.5\%, then after 5 years refinance the remaining balance into a new 25-year loan at 5.2\% APR with \$6,000 in closing costs

Which option minimizes total cost? When does refinancing break even against staying with the original 30-year loan?

% ── Mapping ───────────────────────────────────────────────────────────
\subsection{Real-World Mapping}

\begin{longtable}{ll}
\toprule
\textbf{Real-World Concept} & \textbf{Mathematical Object} \\
\midrule
\endhead
loan balance & $L — principal$ \\
annual percentage rate & $r — monthly rate = APR/12$ \\
loan term & $n — number of monthly payments$ \\
monthly mortgage payment & $PMT — annuity payment$ \\
\bottomrule
\end{longtable}

% ── Constraints ───────────────────────────────────────────────────────
\subsection{Constraints}

\begin{enumerate}
  \item $PMT = L \cdot \frac{r(1+r)^n}{(1+r)^n - 1}$
  \hfill \textit{(annuity payment formula)}
  \item $B_k = B_{k-1}(1+r) - PMT$
  \hfill \textit{(amortization recurrence)}
  \item $B_n = 0$
  \hfill \textit{(loan fully repaid at maturity)}
  \item $\text{Total Interest} = n \cdot PMT - L$
  \hfill \textit{(total interest = payments minus principal)}
\end{enumerate}

% ── Objective / Claim ─────────────────────────────────────────────────
\subsection{Objective}

\begin{equation}
\min_{\text{option}} \; \sum_{k=1}^{n} PMT_k \; + \; \text{closing costs}
\end{equation}

% ── Approach ──────────────────────────────────────────────────────────
\subsection{Approach}

\begin{description}
  \item[Suggested approach:] Annuity PMT + Amortization Schedule + Break-Even Search
  \item[Complexity class:] 
  \item[Available tools:] Annuity PMT + Amortization Schedule + Break-Even Search
\end{description}
% ── Phase B: Solution ────────────────────────────────────────────────
\section{Solution}

\begin{polyaresult}

\textbf{Answer}: cheapest: 15-Year Fixed 5.8\%; refinance break-even at month 26


\textbf{Optimal}: No / Unknown \hfill
\textbf{Feasible}: Yes
\textbf{Algorithm}: Annuity PMT + Amortization Schedule + Break-Even Search \quad $$

\textbf{Time}: 0.002336381992790848s

\textbf{Certificate}: Amortization verified: final balance = \$0.00 for all options.

\end{polyaresult}

% ── Solution Details ──────────────────────────────────────────────────
\subsection{Solution Details}

Home price: \$400,000
Down payment: \$80,000
Loan amount: L = \$320,000

Mortgage Options Comparison:

  30-Year Fixed 6.5\%:
    Monthly payment: PMT = \$2,022.62
    Total interest: \$408,142.38
    Total paid: \$728,143.20

  15-Year Fixed 5.8\%:
    Monthly payment: PMT = \$2,665.89
    Total interest: \$159,859.70
    Total paid: \$479,860.20

  Refinance 30yr->25yr at 5.2\%:
    Monthly payment: PMT = \$1,786.25
    Total interest: \$337,232.57
    Total paid: \$663,232.20

Cheapest option: 15-Year Fixed 5.8\%
Refinance break-even: month 26 (2.2 years)
Refinance net savings: \$64,911

% ── Verification ──────────────────────────────────────────────────────
\subsection{Verification}

No independent verification checks recorded.

% ── Phase C: Interpretation ──────────────────────────────────────────
\section{Interpretation}

% ── The Question & The Answer ─────────────────────────────────────────
\subsection{The Question}
A homebuyer is purchasing a \$400,000 home with an \$80,000 down payment (\$320,000 loan).

\subsection{The Answer}

\begin{polyainsight}[title={Bottom Line}]
Cheapest option: 15-Year Fixed 5.8\%; refinancing saves \$64,911; break-even at month 26.
\end{polyainsight}

% ── What This Means ───────────────────────────────────────────────────
\subsection{What This Means}

- Monthly payments for each option
- Full amortization schedule summaries (total interest, total paid)
- Refinance break-even month (when cumulative savings exceed closing costs)
- Net savings from refinancing over the remaining loan life
- Independent verification that amortization schedules sum correctly

1. PMT calculation: Compute monthly payment using the annuity formula P = L * r(1+r)\textasciicircum{}n / ((1+r)\textasciicircum{}n - 1)
2. Amortization schedule: For each month, split payment into principal and interest, track remaining balance
3. Refinance analysis: After 60 months on the 30-year loan, compute remaining balance, then start a new 25-year schedule at the lower rate plus closing costs
4. Break-even: Find the month where cumulative payment savings from refinancing exceed the closing costs
5. Verify: Independently confirm that principal payments sum to the loan amount and that final balances are zero

% ── Sensitivity Analysis ──────────────────────────────────────────────

% ── Figures ───────────────────────────────────────────────────────────

% ── Recommendations ───────────────────────────────────────────────────
\subsection{Recommendations}

\begin{enumerate}
  \item Compare total interest paid, not just monthly payments, when evaluating mortgage options.
\end{enumerate}

% ── Limitations ───────────────────────────────────────────────────────
\subsection{Limitations}

\begin{polyawarnbox}
\begin{itemize}
  \item Assumes fixed interest rates; adjustable-rate mortgages introduce variability.
  \item Does not account for tax deductions, insurance, or opportunity cost of down payment.
\end{itemize}
\end{polyawarnbox}

% ── Appendix ─────────────────────────────────────────────────────────

\end{document}